% Generated draft: requires mathematical and source review.
\documentclass[11pt]{article}
\usepackage[margin=25mm]{geometry}
\usepackage{fontspec}
\setmainfont{Latin Modern Roman}
\usepackage{amsmath,amssymb,mathtools,unicode-math}
\setmathfont{Latin Modern Math}
\usepackage{enumitem,needspace,xurl,hyperref}
\setlength{\parindent}{0pt}
\setlength{\parskip}{5pt}
\setlength{\emergencystretch}{3em}
\newcommand{\N}{\mathbb{N}}
\newcommand{\Z}{\mathbb{Z}}
\newcommand{\Q}{\mathbb{Q}}
\newcommand{\R}{\mathbb{R}}
\newcommand{\C}{\mathbb{C}}
\newcommand{\F}{\mathbb{F}}
\newcommand{\A}{\mathbb{A}}
\newcommand{\PP}{\mathbb{P}}
\newcommand{\E}{\mathbb{E}}
\newcommand{\eps}{\varepsilon}
\newcommand{\dd}{\,\mathrm{d}}
\begin{document}
\section*{UnsolvedMath ID: 1101105}
\textbf{Source problem code:} catalogue entry
\subsection{Definitions and mathematical statement}
For an integer $n\ge1$ and $A\in\operatorname{GL}(n,\Z)$, let
\[
G_A=\Z^n\rtimes_A\Z,\qquad
(v,k)(w,\ell)=(v+A^k w,k+\ell).
\]
Equivalently, $G_A$ is generated by $\Z^n$ and a stable letter $t$ with $tvt^{-1}=Av$ for $v\in\Z^n$. This is the group mapping torus of the automorphism $A$; it is also the fundamental group of the corresponding torus mapping torus. The zero-dimensional case, if admitted, gives $G_A\cong\Z$ and may be included separately.

Equip each $G_A$ with a word metric from a finite generating set. Two such groups are quasi-isometric if there are a map $f:G_A\to G_B$ and constants $L\ge1$, $C\ge0$ such that
\[
L^{-1}d_A(x,y)-C\le d_B(f(x),f(y))\le Ld_A(x,y)+C
\quad(x,y\in G_A),
\]
and every point of $G_B$ is within distance $C$ of $f(G_A)$. This relation is independent of the chosen finite generating sets.

Bridson's problem asks to classify all $G_A$ up to this equivalence: determine a necessary and sufficient criterion on arbitrary $A\in\operatorname{GL}(n,\Z)$ and $B\in\operatorname{GL}(m,\Z)$, with dimensions allowed to vary, for $G_A$ and $G_B$ to be quasi-isometric. No diagonalizability, spectral restriction, or fixed dimension is imposed by the source. The problem requests a classification; it does not explicitly demand an effective algorithm.

The supplied partial-progress summary organizes matrices, after taking powers, into cases labeled H, N, and M, and reports polynomial-versus-exponential growth separation, rigidity in a ``non-degenerate'' Lie-model case, nilpotent examples including the Heisenberg group, and unresolved cases. However, its referenced Lemmas 1--3, definitions of the cases, and ``honesty caveat'' are absent, and its final sentence is truncated. Consequently no exact case classification is inferred from these labels.

In particular, the summary claims a classification using the multiset of $\log|\lambda_i(A)|$ up to positive rescaling, calls this equivalent to equality of multisets of $|\lambda_i(A)|^p$ and $|\lambda_j(B)|^q$ for positive integers $p,q$, and asserts a corresponding lattice-rigidity statement. Those claims require their omitted hypotheses and precise cited theorems: arbitrary positive rescaling is not by definition rational rescaling, and a classification must address Jordan data and reversal of the cyclic direction. Its assertion that the nilpotent cases contain the entire general nilpotent quasi-isometry classification is also not supported by a reduction in the supplied record. These are unresolved source issues, not adopted classification theorems.
\subsection{Short English statement}
Classify the abelian-by-cyclic groups $\Z^n\rtimes_A\Z$ up to quasi-isometry for all integral invertible matrices and all dimensions. The supplied summary reports partial rigidity and growth results but omits its case definitions and key caveats. Its spectral equivalence and claimed reduction from general nilpotent groups therefore require review and are not treated as established answers.
\subsection{Sources}
\begin{itemize}
\item Local source registry: \url{data/unsolvedmath/problems.json}, permanent UnsolvedMath ID $1101105$, source code catalogue entry, including the supplied PARTIAL-PROGRESS research summary. Dataset attribution: Ulam AI and UnsolvedMath Contributors; source attribution and rights are retained.
\item Mladen Bestvina, \emph{Questions in Geometric Group Theory}, source-list date 2004, Question 11.5, attributed to Bridson; supplied metadata identifies PDF page 19: \url{https://www.math.utah.edu/~bestvina/eprints/questions-updated.pdf}. The record reports manual transcription and July 29, 2026 access, with current-status and rights review flags.
\item The supplied summary names Gromov and Milnor--Wolf for growth distinctions, and Peng I--II, building on work labeled EFW, for Lie-model rigidity. It supplies no full titles, theorem numbers, bibliographic links, definitions of its cases, or the referenced caveat. These references are not independently verified here.
\end{itemize}
\end{document}
